\section{Conclusion}
\label{sec:conclusion}

In this paper we introduced \modelname{}, an inherently interpretable neural
network for tabular classification. High-stakes applications in medicine, finance
and criminal justice require models that are both reasonably accurate and
auditable, and \modelname{} addresses this trade-off geometrically: instead of
unbounded activations such as ReLU, each hidden unit applies a localized
triangular response to a bias-free linear projection. The unit is therefore active
only inside a convex slab bounded by two parallel hyperplanes, and it can be read
as an oblique interval rule with a signed vote per class.

To keep the rule set compact, we combine dual $\ell_1$ regularization with
iterative magnitude pruning. Pruning is a structural mechanism rather than only
compression: zeroing input connections removes features from a rule, and zeroing
an entire row removes the rule. We proved the \emph{exact bias-shifting property
for pruned rules} (Proposition~\ref{prop:bias_shifting}), so a structurally empty
rule can be absorbed into the output bias with zero logit drift; without the
correction, naive removal costs up to 9.01 accuracy points. Rule-adaptive integer
quantization then converts projection weights into small integers, so compact rule
sets can be scored by hand while the output layer remains continuous and
interpretable as signed votes.

The empirical results support a bounded claim. On the single split, the COMPAS
(Global Best) and German Credit models are statistically indistinguishable from
the strongest of 13 baselines. Under the shared folds the German Credit gap is
2.8--3.1 points, and on UCI Adult the models trail the best gradient-boosted
ensemble by 1.52--1.85 points. The interpretable operating point uses 5 rules on
COMPAS, 46 on German Credit and 43 on Adult, with 1.2--2.7 features per rule. On
the 20-class synthetic task the localized-slab representation is much stronger
than tree ensembles, but that benchmark was generated from localized geometry and
should be read as a controlled demonstration rather than as evidence of general
superiority. Integer scorecards retain 95.45--100\% prediction agreement.

Several questions remain open. The architecture targets tabular classification, and
extending convex slabs to regression, survival analysis or unstructured data is
future work. Rules that combine three or more variables require more cognitive
effort than univariate thresholds, and a low rule count alone is not sufficient
for auditability, as the dense German Credit Global Best checkpoint illustrates.
Rule stability is the most important open issue: across five seeds the feature
support is nearly stable (Jaccard $0.91$--$1.00$) but individual rule signatures
are not ($0.02$--$0.13$), and the compact model reduces COMPAS to about five rules
per run with mean rule-set overlap $0.38 \pm 0.15$. A reported rule is therefore
one member of a family of equally accurate rules rather than a uniquely determined
explanation, which makes stability-aware pruning the natural next step. Iterative
pruning also increases offline training time, although deployment cost remains
low. Finally, a transparent model is not automatically a fair model: we do not
evaluate group fairness on COMPAS, and human auditing of proxy attributes remains
essential. Promising directions include causal constraints for recourse,
stability-aware pruning objectives, and compiling quantized scorecards into
low-power edge firmware.
